\subsection{RQ7: How do SLMs compare with substantially larger language models across code generation, error detection, and tool-assisted settings?}

The final research question seeks to understand how the models compare to larger ones in all tasks analyzed previously. All models were compared on the same datasets, with no methodological differences in the comparison, except for the models themselves.

The first analysis concerns the models' performance in code generation, in which models such as Muse Glimmer (80.8\%), GPT-OSS (79.0\%), and Gemma 4 31B (78.7\%) outperformed the others in a two-tailed McNemar test, with a significance threshold of p<0.05, while only four models showed statistically inferior performance. Thus, it is concluded that for this task, some SLMs have comparable or even superior capabilities to larger models. This response demonstrates a leap compared to the study by , because previously, in the same comparison, the difference between the models was clear and favored the larger models.

In comparing error detection capabilities, there is no clear difference between the model sizes. While DeepSeek V4.1 Flash shows the best accuracy, the Muse Glimmer 30B SLM demonstrated superior performance and the best Macro F1 among all evaluated models. Furthermore, most SLMs showed superior performance in several classes, similar to MiMo V2.5. Finally, it is noted that, regardless of the problem size, the main difficulty for both models lies in detecting Runtime Errors, indicating an intrinsic limitation of LLMs in general, even if their detection may be superior to that of static tools.

The feedback from static analysis for hallucination detection also does not generate significant gains in the larger models. The utilization of these models may even be lower than that of SLMs, as the Muse Glimmer 30B, even without the addition of feedback, presented a superior Macro F1 and, with the addition of feedback, still obtained a greater gain in the metric. However, larger models, in general, showed high repair capacity in an iterative cycle, completely reducing code syntax errors and all other types of errors. In contrast, SLMs presented greater difficulty in using the iterative cycle, with some models showing an increase in hallucinations.

It is concluded, therefore, that for this type of task, SLMs such as Muse Glimmer 30B are comparable to larger models, showing similar or even superior performance, as well as similar detection capabilities and tool usage. The main difference lies in some smaller SLMs, such as the Gemma 3 4B, which lags behind not only in performance but also in the use of feedback and the iterative cycle.